\documentclass[11pt]{article}
\usepackage{amsmath}
\usepackage{amssymb}
\usepackage{geometry}
\geometry{margin=1in}
\usepackage[numbers]{natbib}

\title{Vascular Resistance Network Analysis}
\author{}
\date{}

\begin{document}
\maketitle

\section{Vascular Resistance Network Analysis}
\label{sec:resistance_lumping}

\subsection{Skeleton-to-Centerline Graph Construction}

A dense one-dimensional (1-D) centerline graph was constructed directly from the
skeletonized vessel mask.  Every non-zero voxel in the binary skeleton volume was
assigned a unique node positioned at its world-coordinate voxel centre, derived
from the NIfTI affine transformation $\mathbf{A}$:
%
\begin{equation}
  \mathbf{x}_k = \mathbf{A}
  \begin{bmatrix} i_k \\ j_k \\ k_k \\ 1 \end{bmatrix}_{[1:3]},
  \quad k = 1, \ldots, N,
\end{equation}
%
where $(i_k, j_k, k_k)$ are the zero-indexed voxel coordinates of node $k$ and
$N$ is the total number of skeleton voxels.  Undirected edges were placed between
every pair of 26-connected adjacent skeleton voxels, so the mesh topology exactly
mirrors the skeleton connectivity in the original volume.

\subsection{Vessel Radius Estimation}

The local vessel radius at each node was estimated as the Euclidean distance from
that node to the nearest background voxel.  A distance transform was computed on
the binary vessel mask $\Omega$ using the method of
\citet{meijster2000general}:
%
\begin{equation}
  r(\mathbf{x}) = \mathrm{EDT}\!\left(\mathbf{1} - \mathbf{1}_\Omega\right)\!(\mathbf{x}),
\end{equation}
%
where the transform was evaluated with per-axis voxel spacing
$(\Delta x, \Delta y, \Delta z)$ to correctly handle anisotropic acquisitions.
The resulting distance-transform volume was then sampled at each node world
position via nearest-voxel interpolation (using the inverse affine), and values
below $0.1\,\mathrm{mm}$ were clamped to that floor to avoid degenerate
conductances.

\subsection{Hagen-Poiseuille Conductance}

Each edge $e = (i,j)$ was assigned a lumped Hagen-Poiseuille conductance
%
\begin{equation}
  G_{ij} = \frac{\pi \bar{r}_{ij}^{\,4}}{8\,\mu\,L_{ij}},
\end{equation}
%
where $\bar{r}_{ij} = (r_i + r_j)/2$ is the mean nodal radius along the
edge, $L_{ij} = \|\mathbf{x}_i - \mathbf{x}_j\|_2$ is the Euclidean edge
length, and $\mu = 3.5 \times 10^{-3}\,\mathrm{Pa\cdot s}$ is the dynamic
viscosity of blood.  All quantities were converted to SI units before evaluation.

\subsection{Phantom Gap Edges}

Skeletonization and image acquisition artifacts frequently produce multiple
disconnected components.  To obtain a fully connected conductance network,
phantom gap edges were inserted between components separated by less than a
user-specified maximum gap distance $d_\mathrm{max}$.

\paragraph{Candidate endpoint selection.}
For each connected component, candidate gap-bridging nodes were identified as
degree-$\leq 1$ skeleton tips (leaves); for isolated single-node components all
nodes were considered candidates.

\paragraph{MST bridging (default).}
A $k$-d tree was built on all candidate nodes.  All cross-component candidate
pairs within Euclidean distance $d_\mathrm{max}$ were enumerated via
\texttt{query\_pairs}.  For each pair of components $(c_i, c_j)$, only the
closest candidate pair was retained as the representative bridge, giving a
component-level weighted graph.  A minimum spanning tree (MST) was computed on
this component graph to select the minimal set of bridges that ensures full
connectivity.  This guarantees exactly one phantom edge per gap at minimum total
phantom length, avoiding spurious parallel conductance paths that would distort
the flow solution.

\paragraph{k-NN bridging (fallback).}
As an alternative, every cross-component candidate pair within $d_\mathrm{max}$
received a phantom edge.  This matches the behaviour of the legacy MATLAB
implementation but can introduce redundant parallel paths when many tip nodes
are spatially clustered.

\paragraph{Phantom conductance.}
Each phantom edge $(i,j)$ was assigned a penalised conductance
%
\begin{equation}
  G^\mathrm{ph}_{ij} = \frac{\pi \bar{r}_{ij}^{\,4}}{8\,\mu\,\alpha\,L_{ij}},
\end{equation}
%
where $\alpha = 10$ is a dimensionless penalty factor ensuring that phantom
bridges carry negligible flow relative to real vascular connections.

\subsection{Boundary Conditions}

Boundary conditions were assigned automatically from the real-graph topology.
The inlet node was defined as the degree-$\leq 1$ leaf with the largest local
radius (reflecting the largest visible feeding vessel), and prescribed a pressure
$p_\mathrm{in} = 100\,\mathrm{mmHg}$.  All other real-graph leaves were treated
as outlets and prescribed $p_\mathrm{out} = 5\,\mathrm{mmHg}$.  In the absence
of any topological leaves (e.g.\ fully cyclic sub-graphs), the nodes with the
maximum and minimum $z$-coordinates were used as inlet and outlet, respectively.

\subsection{Sparse Conductance System and Solution}

The nodal pressure field $\mathbf{p} \in \mathbb{R}^N$ was obtained by solving
the symmetric sparse Kirchhoff system
%
\begin{equation}
  \mathbf{K}\,\mathbf{p} = \mathbf{q},
\end{equation}
%
where the global conductance (Laplacian) matrix is
%
\begin{equation}
  K_{kl} = \begin{cases}
    \displaystyle\sum_{j \sim k} G_{kj} & k = l, \\[6pt]
    -G_{kl} & k \neq l, \; (k,l) \in \mathcal{E}, \\
    0 & \text{otherwise},
  \end{cases}
\end{equation}
%
and $\mathcal{E}$ is the union of real and phantom edges.

Boundary conditions were incorporated via Schur-complement elimination.  Let
$\mathcal{F}$ and $\mathcal{C}$ denote the sets of free and constrained nodes,
respectively.  The block system
%
\begin{equation}
  \begin{bmatrix} \mathbf{K}_{FF} & \mathbf{K}_{FC} \\
                  \mathbf{K}_{CF} & \mathbf{K}_{CC} \end{bmatrix}
  \begin{bmatrix} \mathbf{p}_F \\ \mathbf{p}_C \end{bmatrix}
  = \mathbf{0}
\end{equation}
%
was reduced to
%
\begin{equation}
  \mathbf{K}_{FF}\,\mathbf{p}_F = -\mathbf{K}_{FC}\,\mathbf{p}_C,
  \label{eq:schur}
\end{equation}
%
which was solved with a direct sparse solver.  Any connected component of the
combined graph that contained no boundary node was pinned to the mean pressure
$p_\mathrm{mean} = (p_\mathrm{in} + p_\mathrm{out})/2$ to prevent a singular
system.  Zero-diagonal rows arising from fully isolated nodes after partitioning
were regularised by setting the diagonal to unity and the right-hand side to
$p_\mathrm{mean}$.

\subsection{Volumetric Flow}

The volumetric flow through each real edge was computed from the pressure
solution as
%
\begin{equation}
  Q_{ij} = G_{ij}\,(p_i - p_j),
\end{equation}
%
with results reported in $\mathrm{mm^3\,s^{-1}}$.  The full pressure, radius,
flow, conductance, and edge-length fields were written back to the centerline
mesh in VTK PolyData format (VTP) for visualisation in ParaView.

\bibliographystyle{plainnat}
\bibliography{references}

\end{document}
